\section*{Chimie (7 points)}
La nature des systèmes chimiques dépend des corps mis en jeu. Le suivi de l'évolution de ces systèmes peut se faire par différentes méthodes, physiques ou chimiques, ce qui amène à faire des mesures et à déterminer des grandeurs caractéristiques de ces systèmes et des transformations auxquelles ils sont soumis.\\
Cet exercice vise l'étude de certains systèmes faisant intervenir l'acide méthanoïque.

\subsection{Étude de la solution d'acide méthanoïque}
On considère une solution aqueuse $\mathrm{(S_{1})}$ d'acide méthanoïque
$\ce{HCOOH}$ de concentration molaire ${C_{1}=\qty{5e-3}{\mol\per\L}}$. La
mesure de la conductivité de cette solution, à $\qty{25}{\degreeCelsius}$, donne
$\sigma_{1}=\qty{33}{\mS\per\m}$.

\begin{donnees}
    \begin{itemize}
        \item Conductivités molaires ioniques~:

              $\lambda_{1}=\lambda(\ce{H3O^+})=\qty{35.0}{\mS\square\m\per\mol}$~;
              $\lambda_{2}=\lambda(\ce{HCOO^-})=\qty{5.5}{\mS\square\m\per\mol}$

        \item L'effet des ions $\ce{HO^{-}_{(aq)}}$ sur la conductivité de la
              solution est supposé négligeable.
        \item La conductivité $\sigma$ d'une solution s'écrit en fonction des
              concentrations molaires effectives des ions $\ce{X_{i}}$ et de
              leurs conductivités molaires ioniques $\lambda_{i}$ comme suit~:
              $\sigma=\sum\lambda_{i}\bigl[\ce{X_{i}}\bigr]$.
    \end{itemize}
\end{donnees}
\begin{enumerate}
    \item Écrire l'équation chimique modélisant la réaction entre l'acide
          méthanoïque et l'eau.
    \item Montrer que
          $\bigl[\ce{H3O^{+}_{(aq)}}\bigr]=\qty{8.15e-4}{\mol\per\L}$.
    \item Calculer le taux d'avancement final $\tau_{1}$ de la réaction.
          Conclure.
    \item Montrer que la valeur du quotient de réaction à l'état d'équilibre du
          système chimique est ${Q_{r_{1},\eq}=\num{1.59e-4}}$.
    \item On dilue à $\qty{25}{\degreeCelsius}$ la solution $\mathrm{(S_{1})}$,
          pour obtenir une solution aqueuse $\mathrm{(S_{2})}$ de concentration
          molaire $C_{2}$.
    
          Donner, en justifiant, la valeur du quotient de réaction
          $Q_{r_{2},\eq}$ du système chimique à l'état d'équilibre.
\end{enumerate}




\subsection{Exploitation du critère d'évolution}
On considère le système chimique obtenu en mélangeant les quantités de matière
suivantes~: ${n_{1}=\qty{1.5e-2}{\mol}}$ d'acide nitreux $\ce{HNO2}$,
$n_{2}=\qty{3e-2}{\mol}$ de méthanoate de sodium $\ce{Na^+_{(aq)} +
HCOO^-_{(aq)}}$, $n_{3}=\qty{3e-2}{\mol}$ de nitrite de sodium $\ce{Na^+_{(aq)}
+ NO^-_2_{(aq)}}$ et $n_{4}=\qty{1.5e-2}{\mol}$ d'acide méthanoïque
$\ce{HCOOH}$. Soit $V$ le volume total du mélange réactionnel.

L'équation de la réaction entre l'acide nitreux $\ce{HNO2}$ et les ions
méthanoate $\ce{HCOO^-_{(aq)}}$ s'écrit~:
\[
    \ce{HNO2_{(aq)} + HCOO^-_{(aq)} <=>[(1)][(2)] NO^-_2_{(aq)} + HCOOH_{(aq)}}
\]

\begin{donnees}
    $\mathrm{pK_{A_{1}}}=
    \mathrm{pK_{A}}(\ce{HNO2_{(aq)}}\ttslash{}\ce{NO^-_2_{(aq)}})=3,2$~; \quad
    $\mathrm{pK_{A_2}}=
    \mathrm{pK_{A}}(\ce{HCOOH_{(aq)}}\ttslash{}\ce{HCOO^-_{(aq)}})=3,8$
\end{donnees}

\begin{enumerate}
    \item Déterminer la valeur du quotient de réaction $Q_{r,i}$ à l'état
          initial du système chimique.
    \item Montrer que la constante d'équilibre $\mathrm{K}$ associée à
          l'équation chimique précédente s'écrit~:
          ${\mathrm{K}=10^{(\mathrm{pK_{A_2}-pK_{A_1}})}}$.
          Calculer la valeur de $\mathrm{K}$.
    \item Indiquer, en justifiant, dans quel sens évolue spontanément le système
          chimique à partir de son état initial.
\end{enumerate}




\section*{3. Suivi temporel d'une réaction chimique}
On prépare un mélange contenant initialement de l'acide méthanoïque et un alcool. Le suivi de la réaction a permis de tracer la courbe ci-contre représentant l'avancement $x$ de la réaction en fonction du temps.\\
Le volume total du mélange est $V=88 m L$.\\
3.1. Déterminer graphiquement :\\
a. la valeur de l'avancement final $x_{f}$ de la réaction.\\
b. la valeur du temps de demi-réaction $t_{1 / 2}$.\\
c. la valeur de la vitesse volumique de la réaction à l'instant $t_{0}=0$ en unité ( $m o l \cdot L^{-1} \cdot h^{-1}$ ).\\
\includegraphics[max width=\textwidth, alt={}, center]{3722cc87-920c-4836-a0f1-9718f4bdb853-3_607_832_303_1135}\\
3.2. Interpréter qualitativement la variation de la vitesse volumique de la réaction.

